\documentclass{ctexart}
\usepackage{template}
\usepackage{esint}

\begin{document}\pagestyle{empty}
\begin{center}\Large
    北京大学2025-2026学年第二学期期中试题\\
    高等数学B2
\end{center}
\begin{problem}[1.(10\songti{分})]
设$D$是平面上的闭圆盘$x^2+y^2\leqslant2x+2y$.计算二重积分
    \[\iint_Dxy\di x\di y.\]
\end{problem}
\begin{solution}
令$u=x-1$, $v=y-1$,则积分区域化为
\[D_0=\{(u,v):u^2+v^2\leqslant2\},\qquad xy=(u+1)(v+1)=uv+u+v+1.\]
由于$D_0$分别关于$u$轴和$v$轴对称,其中$uv$, $u$, $v$的积分均为$0$.于是
\[\iint_Dxy\di x\di y
=\iint_{D_0}(uv+u+v+1)\di u\di v
=\iint_{D_0}1\di u\di v=2\pi.\]
\end{solution}
\begin{problem}[2.(15\songti{分})]
设$\Omega$是$\R^3$中球体$x^2+y^2+z^2\leqslant R^2$ $(R>0)$与锥体$z\geqslant\sqrt{x^2+y^2}$的交集.计算三重积分
    \[\iiint_\Omega(x^2+y^2)\di x\di y\di z.\]
\end{problem}
\begin{solution}
采用球坐标
\[x=r\sin\varphi\cos\theta,\quad y=r\sin\varphi\sin\theta,\quad z=r\cos\varphi.\]
锥体条件等价于$0\leqslant\varphi\leqslant\pi/4$,因而积分区域为
\[0\leqslant r\leqslant R,\quad0\leqslant\varphi\leqslant\frac\pi4,\quad0\leqslant\theta\leqslant2\pi.\]
由$x^2+y^2=r^2\sin^2\varphi$及体积元$r^2\sin\varphi\di r\di\varphi\di\theta$,得到
\[\begin{aligned}
\iiint_\Omega(x^2+y^2)\di x\di y\di z
&=\int_0^{2\pi}\di\theta\int_0^{\pi/4}\sin^3\varphi\di\varphi
\int_0^Rr^4\di r\\
&=\frac{2\pi R^5}{5}\left[-\cos\varphi+\frac{\cos^3\varphi}{3}\right]_0^{\pi/4}\\
&=\frac{2\pi R^5}{5}\left(\frac23-\frac{5\sqrt2}{12}\right)
=\frac{\pi R^5}{30}(8-5\sqrt2).
\end{aligned}\]
\end{solution}
\begin{problem}[3.(15\songti{分})]
设$C(R)$是平面上的圆周$x^2+y^2=R^2$ $(R>0)$,定向为逆时针方向.记第二型曲线积分
    \[I(R)=\oint_{C(R)}\dfrac{x\di y-y\di x}{(x^2+xy+y^2)^2}.\]
    证明
    \[\dfrac{13\pi}{9R^2}\leqslant I(R)\leqslant\dfrac{5\pi}{R^2}.\]
\end{problem}
\begin{solution}
取圆周的逆时针参数方程$x=R\cos\theta$, $y=R\sin\theta$, $0\leqslant\theta\leqslant2\pi$.于是
\[x\di y-y\di x=R^2\di\theta,\qquad
x^2+xy+y^2=R^2\left(1+\frac12\sin2\theta\right),\]
从而
\[I(R)=\frac1{R^2}\int_0^{2\pi}\frac{\di\theta}{\left(1+\frac12\sin2\theta\right)^2}.\]
令$E_+=\{\theta\in[0,2\pi]:\sin2\theta\geqslant0\}$,
$E_-=\{\theta\in[0,2\pi]:\sin2\theta<0\}$.这两个集合的总长度均为$\pi$.
在$E_+$上有$1\leqslant1+\frac12\sin2\theta\leqslant\frac32$,在$E_-$上有
$\frac12\leqslant1+\frac12\sin2\theta\leqslant1$.因此
\[\frac49\leqslant\frac1{\left(1+\frac12\sin2\theta\right)^2}\leqslant1
\quad(\theta\in E_+),\]
\[1\leqslant\frac1{\left(1+\frac12\sin2\theta\right)^2}\leqslant4
\quad(\theta\in E_-).\]
分别积分后相加,可得
\[\frac1{R^2}\left(\frac49\pi+\pi\right)
\leqslant I(R)\leqslant\frac1{R^2}(\pi+4\pi),\]
即$\displaystyle\frac{13\pi}{9R^2}\leqslant I(R)\leqslant\frac{5\pi}{R^2}$.
\end{solution}
\begin{problem}[4.(15\songti{分})]
设$\Sigma$是$\R^3$中圆柱面$x^2+z^2=R^2$ $(R>0)$上满足$x^2+y^2\leqslant R^2$, $z\geqslant0$的部分.计算第一型曲面积分
    \[\iint_\Sigma z^2\di S.\]
\end{problem}
\begin{solution}
取曲面的参数方程
\[\boldsymbol r(\theta,y)=(R\cos\theta,y,R\sin\theta),\qquad0\leqslant\theta\leqslant\pi.\]
由$x^2+y^2\leqslant R^2$可知$-R\sin\theta\leqslant y\leqslant R\sin\theta$.又有
\[\left|\frac{\partial\boldsymbol r}{\partial\theta}\times
\frac{\partial\boldsymbol r}{\partial y}\right|=R,\qquad z^2=R^2\sin^2\theta.\]
因此
\[\begin{aligned}
\iint_\Sigma z^2\di S
&=\int_0^\pi\int_{-R\sin\theta}^{R\sin\theta}R^3\sin^2\theta\di y\di\theta\\
&=2R^4\int_0^\pi\sin^3\theta\di\theta
=2R^4\left[-\cos\theta+\frac{\cos^3\theta}{3}\right]_0^\pi
=\frac83R^4.
\end{aligned}\]
\end{solution}
\begin{problem}[5.(15\songti{分})]
设$S$是$\R^3$中圆柱面$(x-1)^2+y^2=1$ $(0\leqslant z\leqslant h)$的外侧.计算第二型曲面积分
    \[I=\iint_Sx^2\di y\di z+y^2\di z\di x+z^2\di x\di y.\]
\end{problem}
\begin{solution}
当$h>0$时,令
\[V=\{(x,y,z):(x-1)^2+y^2\leqslant1,\ 0\leqslant z\leqslant h\}.\]
用底面$S_0$和顶面$S_h$将$S$补成闭曲面,均取外侧.对向量场
$\boldsymbol F=(x^2,y^2,z^2)$使用Gauss公式,得到
\[\iint_{\partial V}\boldsymbol F\cdot\boldsymbol n\di S
=\iiint_V(2x+2y+2z)\di V.\]
横截面是以$(1,0)$为圆心的单位圆盘,其面积为$\pi$,由对称性有
\[\iiint_Vx\di V=\pi h,\qquad\iiint_Vy\di V=0,\qquad
\iiint_Vz\di V=\pi\int_0^hz\di z=\frac{\pi h^2}{2}.\]
因此闭曲面的外向通量为$2\pi h+\pi h^2$.底面上$z=0$,其通量为$0$;
顶面的外法向为$(0,0,1)$,其通量为$\pi h^2$.于是侧面的积分为
\[I=2\pi h+\pi h^2-0-\pi h^2=2\pi h.\]
$h=0$时曲面退化,积分为$0$,仍满足此式.若将$h<0$也纳入讨论,则题设集合为空,积分为$0$.
\end{solution}
\begin{problem}[6.(15\songti{分})]
求微分方程
    \[y''+4y=\sin x+\cos2x\]
    的通解.
\end{problem}
\begin{solution}
先求对应齐次方程$y''+4y=0$的通解.特征方程为$\lambda^2+4=0$,因此
\[y_h=C_1\cos2x+C_2\sin2x.\]
对非齐次项$\sin x$,设特解为$y_1=A\sin x$,代入得到$3A\sin x=\sin x$,故$A=1/3$.
对非齐次项$\cos2x$,由于发生共振,设特解为$y_2=Bx\sin2x$.计算可得
\[y_2''+4y_2=4B\cos2x,\]
从而$B=1/4$.根据线性方程解的叠加原理,原方程的通解为
\[y=C_1\cos2x+C_2\sin2x+\frac13\sin x+\frac{x}{4}\sin2x,
\qquad C_1,C_2\in\R.\]
\end{solution}
\begin{problem}[7.(15\songti{分})]
求满足方程
    \[y'(x)=xy(x)+x\int_0^1ty(t)\di t\]
    的所有可导函数$y(x)$, $x\in[0,1]$.
\end{problem}
\begin{solution}
令$\displaystyle A=\int_0^1ty(t)\di t$.这里$t$是积分变量,$A$是与$x$无关的常数.
原方程化为一阶线性方程
\[y'(x)-xy(x)=Ax.\]
乘以积分因子$\e^{-x^2/2}$,有
\[\left(\e^{-x^2/2}y(x)\right)'=Ax\e^{-x^2/2}
=-A\left(\e^{-x^2/2}\right)'.\]
于是存在常数$C$,使得
\[y(x)=C\e^{x^2/2}-A.\]
再将此式代回$A$的定义,得到
\[A=C\int_0^1t\e^{t^2/2}\di t-A\int_0^1t\di t
=C(\sqrt\e-1)-\frac A2.\]
故$\displaystyle A=\frac23C(\sqrt\e-1)$,从而所有可能的解均具有形式
\[y(x)=C\left(\e^{x^2/2}-\frac23(\sqrt\e-1)\right),
\qquad0\leqslant x\leqslant1,\quad C\in\R.\]
反过来,对任意$C\in\R$,此函数可导,且其积分$\int_0^1ty(t)\di t$恰为
$\frac23C(\sqrt\e-1)$.因此$xy(x)+xA=Cx\e^{x^2/2}=y'(x)$,
确实满足原方程.上述函数即为全部解.
\end{solution}
\end{document}
